\documentclass{report}
\usepackage{amsmath,amssymb}
\setlength{\parindent}{0mm}
\setlength{\parskip}{1em}
\begin{document}
\begin{center}
\rule{6in}{1pt} \
{\large Kai Hertel \\
{\bf An iterative solver combined with extrapolation for steady-state Maxwell problems}}

Universitaet Erlangen-Nuernberg \\ Lehrstuhl fuer Systemsimulation \\ Cauerstr 11 \\ 91058 Erlangen \\ Germany
\\
{\tt kai.hertel@cs.fau.de}\\
Christoph Pflaum\end{center}

We present an iterative finite difference frequency domain (FDFD) solver
for a large scale time-harmonic Maxwell problem. The solver is based on
an explicit leap-frog power iteration and is enhanced by a Prony method
extrapolation scheme to speed up convergence.
We take a look at the algebraic formulation of the problem, iteration and
extrapolation schemes, and discuss challenges and results.
The primary application is a thin-film solar cell simulation that
consists of both layers of positive and negative permittivities with
textured interfaces in between.
Electrical permittivity parameters of the layer media as well as the
textured interface profiles are important model parameters.
Resulting quantities of interest are the simulated solar cell's
absorption behavior, specifically its short circuit current density and
its characteristic frequency dependent quantum efficiency.
Considered boundary conditions are periodic or of Dirichlet type in
connection with a perfectly matched layer (PML) to simulate free space
propagation.
The periodicity of the domain with respect to textured interfaces can be
enforced either by a reformulation of the periodic boundary condition or
a transition layer that extends on the classical approach to achieving
periodicity by means of symmetrical domain extensions.
The discussed method scales to large problem sizes, and obtained results
are in good agreement with physical measurements of manufactured
structures.


\end{document}
